\documentclass{article}

\usepackage{graphicx} % Required for inserting images
\usepackage{authblk} % For affiliations
\usepackage{enumerate} % For ordered lists
\usepackage{parskip} % Adds white space between paragraphs
\usepackage{orcidlink} % Insert ORCID links
\usepackage{abstract} % Customise abstract
\usepackage{xurl} % To add URLs (which will line break)
\usepackage{hyperref} % To add URLs
\usepackage{geometry} % Tables
\usepackage{booktabs} % Tables (styling)
\usepackage{tabulary,lipsum} % Tables (e.g. J/C/L/R column types)
\usepackage{float} % Tables (prevent repositioning)
\usepackage{enumitem} % Use alternatives to 123 in ordered list
\usepackage{forest} % To create directory tree
\usepackage{xcolor} % For highlights and shading
\usepackage{soul} % For highlights
\usepackage{framed} % To make shaded abstract
\usepackage{enumitem,amssymb} % To make check box lists
\usepackage{multicol} % Columns
\usepackage{titlesec} % To make alternative section titles
\usepackage{clock} % Create clock symbol
\usepackage[clock]{ifsym} % Create clock symbol
\usepackage{amsmath} % Create tick symbol
\usepackage{makecell} % To add forced line breaks within table cells
\usepackage{changepage} % For adjustwidth

% Checkboxes
\newlist{todolist}{itemize}{2}
\setlist[todolist]{label=$\square$}

% Change font
\usepackage{cmbright}

% References using biblatex
\usepackage[
    sorting=none,
    autocite=superscript
]{biblatex}
\bibliography{references}

% Set depth of table of contents
\setcounter{tocdepth}{2}

% Blue highlight
\definecolor{highlightblue}{HTML}{e8f2f6}
\DeclareRobustCommand{\hlblue}[1]{{\sethlcolor{highlightblue}\hl{#1}}}
    
% Create shaded abstract
% ---------------------------------------------------
% Credit to JamesT (https://tex.stackexchange.com/questions/675419/text-is-not-centered-with-colorbox)
\definecolor{abstractcolour}{HTML}{D1E6EE}
\colorlet{shadecolor}{abstractcolour}
\renewenvironment{abstract} 
     {
    \vfill%
      \list{}{
        \setlength{\leftmargin}{.3cm}%
        \setlength{\rightmargin}{\leftmargin}%
      \vfill%
      }%
      \item\relax}
 {\endlist}
\setlength{\parindent}{0mm} 
% ---------------------------------------------------

% Create shaded headings
% ---------------------------------------------------
% Logo is #1982AB. Chose paler colous for the heading shading.
\definecolor{section}{HTML}{A3CDDD}
\newcommand{\colorsection}[1]{\colorbox{section}{\parbox{\dimexpr\textwidth-11\fboxsep}{#1}}}
\titleformat
    {name=\section} % Command
    [block] % Shape
    {\normalfont\Large\bfseries} % Format
    {\thesection} % Label
    {1em} % Space between number and title
    {\colorsection} % Before code
    [] % After code

\definecolor{subsection}{HTML}{A3CDDD}
\newcommand{\colorsubsection}[1]{\colorbox{subsection}{\parbox{\dimexpr\textwidth-13\fboxsep}{#1}}}
\titleformat{name=\subsection}[block]
    {\normalfont\large\bfseries}
    {\thesubsection}
    {1em}
    {\colorsubsection}

\definecolor{subsubsection}{HTML}{D1E6EE}
\newcommand{\colorsubsubsection}[1]{\colorbox{subsubsection}{\parbox{\dimexpr\textwidth-14\fboxsep}{#1}}}
\titleformat{name=\subsubsection}[block]
    {\normalfont\normalsize\bfseries}
    {\thesubsubsection}
    {1em}
    {\colorsubsubsection}
% ---------------------------------------------------

\begin{document}

% Title page and TOC
\title{
    % Move higher on page
    % \vspace{-3cm}
    % Include image then 1cm gap
    \includegraphics[width=9cm]{images/stars_logo_blue_text.png}\\[1cm]
    % Title
    \textbf{Protocol for assessing the computational reproducibility of simulation models on STARS} \\
    % Sub-title
    \large Zenodo pre-registration
}

% Authors
\author[1]{\orcidlink{0000-0002-6596-3479} Amy Heather}
\author[1]{\orcidlink{0000-0003-2631-4481} Thomas Monks}
\author[2]{\orcidlink{0000-0001-5274-5037} Alison Harper}
% Include line break (to prevent name breaking over two lines)
\author[2]{\\ \orcidlink{0000-0002-2204-8924} Navonil Mustafee}
\author[3]{\orcidlink{0000-0003-1263-2286} Andy Mayne}

% Affiliations (ensure update with numbers used above)
\affil[1]{\footnotesize University of Exeter Medical School, Exeter, UK}
\affil[2]{\footnotesize University of Exeter Business School, Exeter, UK}
\affil[3]{\footnotesize Taunton and Somerset NHS Foundation Trust, UK}

% No date
\date{}

\maketitle

\textbf{Date of submission:} TBC

% Add space before abstract
\vspace{0.5cm}

% Create abstract in a shaded box
\begin{shaded}
    \begin{abstract}
        This protocol outlines how we plan to reuse available artifacts to reproduce results from published simulation studies. This forms part of the project STARS: "Sharing Tools and Artefacts for Reproducible Simulations in healthcare". It will be utilised to assessed the computational reproducibility of published and newly developed simulation models in Python and R.
    \end{abstract}
\end{shaded}

% Add space after abstract
\vspace{0.5cm}

\newpage
\tableofcontents

% Main body of protocol
\section{Introduction}

In this protocol, we are focused on the \textbf{computational reproducibility} of simulation models. This is defined as the ability to get consistent results with a prior study when using  the same data and methods as that study. We are focusing on models developed using \textbf{Python and R}, as these are popular free and open-source software (FOSS) for the development of models like discrete-event simulation (DES).\autocite{monks_computer_2023}

The protocol will be used to assess the computational reproducibility of \textbf{six DES models}. These are selected from the models identified by \textbf{Monks and Harper 2023}.\autocite{monks_computer_2023} The selection criteria are that: (i) the model code is publicly available, (ii) the model is created using Python or R, and (iii) the code has an open license (either already published, or added upon request from the STARS team).

Throughout the study, results will be openly available and shared via a \textbf{Quarto website}. This will compile information on the reproduction of the article. This includes the notebooks (.ipynb or .Rmd) producing the items in the scope, as well as a chronological log of work using Quarto blog posts, and then later, the reproduction report and detailed study results.

The protocol will often refer to our \textbf{template repository} which can be viewed here -\url{https://github.com/pythonhealthdatascience/stars_reproduction_template}. As you will see, one of the first steps for researchers will be to use this template to set-up their repository for the reproducibility assessment. \hl{CHANGE TO ZENODO CITATION ALONGSIDE URL BEFORE PUBLISHING}

There are \textbf{three key stages} to this protocol which you should \textbf{work through in the following order}:
\begin{enumerate}
    \item Assessment of computational reproducibility
    \item Evaluation against guidelines
    \item Research compendium and report
\end{enumerate}

However, we first introduce two important processes that will need to be \textbf{completed during some or all of the stages}:
\begin{itemize}
    \item Keeping a detailed record of work using a \textbf{logbook}
    \item \textbf{Timing} how long tasks take to complete
\end{itemize}

\newpage
\section{Logbook}

Throughout \textbf{all of the stages} below, you should keep a \textbf{logbook}. This is a detailed record of work recorded using \textbf{Quarto blog posts}, using the template provided. As suggested by Ayllón et al. 2021\autocite{ayllon_keeping_2021} in their guidelines for keeping modelling notebooks, these posts will be \textbf{daily}, dated, chronological entries. \textbf{Tags} will be used to help indicate the activity on each day, and enable posts to be filtered by activity. Keeping a detailed log will support later understanding of what was done, and support preparing of final documents like the reproduction report.

Each entry in the logbook should contain the:
\begin{itemize}
    \item \textbf{Researcher name} and \textbf{date}
    \item \textbf{Tags} (e.g. setup, scope, read, reproduce, guidelines, compendium, report)
    \item \textbf{Time} spent on tasks (if applicable)
    \item \textbf{Comprehensive} record of work. This should include record of working through each stage in the protocol, detailing \textbf{successes} and any \textbf{issues} faced, any \textbf{solutions} found to problems, and any \textbf{changes made to the model code} (noting where and how the code was changed). It may be relevant to include links to particular versions of a file or repository such as via the Git commit history.
    \item Clear statement if and when each item in the scope is considered to have been \textbf{successfully reproduced}.
\end{itemize}

\vspace{0.5cm}
\section{Timing}

During the first stage of your study (\textbf{assessment of computational reproducibility}) you should \textbf{time how long each task takes}. These times should be recorded within the logbook alongside each activity (e.g. 12:10 to 12:45). The times should be monitored with a \textbf{maximum of 40 hours} allowed for attempting to reproduce the study, as in Krafczyk et al. 2021.\autocite{krafczyk_learning_2021} This cut-off is implemented as we anticipate there would be little more to learn from spending longer than that time on reproducing a single study. The timings are just implemented for the first stage, as our focus in timing this work is to understand (and ultimately also limit) how long we spend reproducing each study.

The only exceptions to this timing are:
\begin{itemize}
    \item \textbf{Computation time}. This is at the researcher's discretion. For example, you should include short run times where you are still continuously working on the study. You should exclude longer run times where you are no longer working towards the assessment of computational reproducibility - for example, if you set the simulation to run for five minutes whilst going to make a cup of tea.
    \item \textbf{Consensus on the scope}. You should record the time spent by the primary researcher in reading the article and identifying what they think is in scope. However, you should not include the time spent discussing the scope with other researchers on the team, nor any of the time that they spent reading the article or thinking about scope.
\end{itemize}
\section{First stage: reproduction} \label{sec:reproduce}

\textbf{Remember!} Record progress and the \textbf{time} spent on tasks in the \textbf{logbook} (as described in section \ref{sec:timing}).

\vspace{0.5cm}
\subsection{Set-up}

\subsubsection{Inform the authors}

This step may have already been completed  (for example, if the author has already been contacted to request the addition of an open license to their repository). If not, the researcher should:
\begin{enumerate}
    \item \textbf{Email the corresponding author} about the study, including a \textbf{copy of the study protocol}. A \textbf{template} for this email is provided in Appendix \ref{appendix:email}.
    \item If the email rebounds, search online for a \textbf{recent email address} for any of the study authors.
\end{enumerate}

If the study authors do not reply, there is no need to follow up on this email asking for a response, as the email simply serves to notify the study authors and does not require a response.

\vspace{0.5cm}
\subsubsection{Create repository using template}

To set up the repository:

\begin{enumerate}
    \item Go to \url{https://github.com/pythonhealthdatascience/stars_reproduction_template} and select the ``\textbf{Use this template}" button.
    \item \textbf{Set up the repository} in 'pythonhealthdatascience' with the name 'stars-reproduce-surname-year' and description 'Assessing the computational reproducibility of surname et al. year as part of STARS.' Here, 'surname' and 'year' refer to the article being reproduced.
    \item \textbf{Update the template} so that it is specific to the study being evaluated (e.g. titles, links, CITATION.cff).
    \item Set up the provided \textbf{Python environment} ('requirements.txt') for creation of the Quarto site.
\end{enumerate}

\vspace{0.5cm}
\subsubsection{Upload code to the repository}

We require that the model code has been shared under an open license. Hence, the researcher should be free to upload this to the repository. They should make a full copy of their code repository, uploading any \textbf{code} and related artefacts.

\vspace{0.5cm}
\subsubsection{Update license for repository if required}

Check which license the original study used for their code. By default, the template includes an MIT license, but \textbf{this may need to be changed if the authors used a different license}, so it is compatible.

Following the guidance of the Turing Way\autocite{the_turing_way_community_turing_2022} and R packages book,\autocite{wickham_12_2023} the license does not need to be modified according to the packages used, unless code from that package is embedded within the work (for example, copy+and+paste a function) or if it is distributed as a binary with the work (i.e. bundled and stored with the work, rather than setting it to be exported from somewhere like PyPI or CRAN). If the code simply imports and uses functions from packages, this is not assumed to be derivative work, and hence a permissive license can be chosen. Likewise, the license can be permissive when using Docker\autocite{merkel_docker_2014} as we do not distribute the Docker image itself, but instead distribute the Dockerfile or refer to someone else to distribute (such as the GitHub container registry).\autocite{the_linux_foundation_docker_nodate}

\vspace{0.5cm}
\subsubsection{Upload journal article to the repository}

Before uploading the article to the repository, check which \textbf{license the article has been distributed under}. This is likely to be in a ``copyright" section of the article. If it is distributed under a permissive license (e.g. CC-BY), then the researcher should be free to upload the article and artefacts to the repository. However, if the article does not have a license enabling reuse in this manner (which is likely if it is restricted access/behind a paywall), then they will likely be unable to. Although permission can be requested, this will often incur costs.

In these cases, the researcher should search for a \textbf{green open-access version} of the article. For example, copies of the article available on a pre-print site or in an institutional repository. If a green open-access article is identified, it can be uploaded to the repository - although it should be compared against the paywalled article to check for any differences between them before upload.

If an openly licensed version of the article cannot be located then \textbf{links to the original work} should be provided instead of uploading the article itself. If located though, the researcher should \textbf{upload all available materials} from the article to the 'original\_study/' folder. This includes:
\begin{itemize}
    \item The \textbf{journal article} and any \textbf{supplementary materials}
    \item Each \textbf{table and figure} from the main journal article as individual digital objects (e.g. .jpeg, .png) (not including those from the appendices/supplementary materials). These are uploaded to the repository separately as we will later use the files when describing the scope of the study on the Quarto site, and we will compare our results against them if they are within scope.
\end{itemize}

The researcher should amend the provided \textbf{template page} ('quarto\_site/study\_publication.qmd') to display the PDFs or links to the journal article, and provide a link to the code. If using a green open-access article, a link to the paywalled article should also be included.

\vspace{0.5cm}
\subsection{Scope of reproduction}

\subsubsection{Read the journal article}

\textbf{Read through the journal article} (but not yet looking into the code or data). Any notes can be recorded within the logbook. These can be within a collapsible callout to aid readability of the log.

\vspace{0.5cm}
\subsubsection{Define scope of reproduction}

The next step is to define the scope of the reproducibility study - in other words, what parts of the paper we intend to reproduce. This should be focused on the \textbf{results of the simulation} (rather than other results like description of the sample). To identify the scope:

\begin{enumerate}
    \item Look through each of the \textbf{tables and figures} in the article (excluding the supplementary material) and identify whether they are within scope or not (i.e. do they present results of the simulation).
    \item Identify any \textbf{'key results' in the text} of the article that are within scope. These are results that are highlighted within the \textbf{abstract or results} section. This may include items from the supplementary materials if referred to in the text.
    \item Evaluate whether the identified 'key results' are \textbf{already covered by the tables and figures} from the article. If not, they should be included in the scope.
    \item Make a \textbf{consensus decision on scope with at least one other team member}.
\end{enumerate}

Notes from thoughts and discussions around the potential scope can be recorded within the \textbf{logbook}.

The criteria for what is considered part of the scope is adapted from Wood et al. (2018).\autocite{wood_replication_2018, wood_push_2018}

\vspace{0.5cm}
\subsubsection{Compile items in scope}

Once the scope has been decided, upload each item to the 'original\_study/' folder. If the article does not have permissions to enable this, then the researcher can upload these materials to a separate \textbf{private repository} for their own reference, but cannot display these within the public repository.

\begin{itemize}
    \item For \textbf{tables}, download a \textbf{CSV} version of each if available (or convert into CSV, if not).
    \item For \textbf{figures}, download the \textbf{highest-quality} version of each figure that is available (if not already saved in the repository).
    \item For results described in the \textbf{text} (but not captured in a table or figure), record in a format appropriate to then later compare against (for example, within a \textbf{CSV}).
\end{itemize}

The researcher should amend the provided \textbf{template page} ('evaluation/scope.qmd') to display each of the items.

\vspace{0.5cm}
\subsubsection{Archive scope on Zenodo}

With the organisation linked to Zenodo, \textbf{toggle Zenodo to preserve} that repository, and then create a \textbf{release on GitHub}, which Zenodo will then automatically download and register with a DOI. The release should be recorded within the \textbf{changelog}.

This release serves as a public registration of the intended scope of the reproduction, archiving the repository at this point (prior to having started using or really looking at the code). As stated in the ACRe guide,\autocite{berkeley_initiative_for_transparency_in_the_social_sciences_guide_2022} it is important that the scope is defined at the start of the study, and publicly archived so as not to be amended during the course of the study.\autocite{berkeley_initiative_for_transparency_in_the_social_sciences_guide_2022}

\vspace{0.5cm}
\subsection{Familiarise with artefacts}

\subsubsection{Look over the code/data}

\textbf{Browse through any code and data} uploaded to the repository. The aim of this step is for the researcher to familiarise with the materials before setting up the environment or running the code. There are no requirements for how they should do this, but some suggestions include:
\begin{itemize}
    \item In the logbook, record a one-sentence description of each file and tree of the uploaded files (as suggested by Ayllón et al. (2021)\autocite{ayllon_keeping_2021})
    \item Look for sections of code that produce items within the scope and record this within the logbook (as in Krafczyk et al. 2021\autocite{krafczyk_learning_2021}).
\end{itemize}

\vspace{0.5cm}
\subsection{Set up environment}

Identify the \textbf{software packages and their versions}, as used by the original study authors. This may be provided in an \textbf{environment file or within the article}. If not, the researcher should:
\begin{itemize}
    \item Identify packages from those named to \textbf{import} within each of the scripts
    \item Select versions by looking at the \textbf{version history} for that package on a package repository (e.g. PyPI, CRAN), and identifying a version whose release date is \textbf{closest to but still prior to} the date of the code archive or paper publication (whichever is earliest).
\end{itemize}

Use or set up an environment file with the identified dependencies, and \textbf{create the environment}. Researchers should use simple methods for environment management such as Conda or VirtualEnv in Python, and renv in R. For simplicity, we are not requiring that they match the operating system used.

\textbf{Important!} This should be set up within the '\textbf{reproduction/}' folder, whilst the 'original\_study/' folder should remain untouched.

\vspace{0.5cm}
\subsection{Attempt to reproduce items in scope}

This stage is an \textbf{iterative} process of running the code and attempting to reproduce the items in the scope. For this stage, the researcher should:
\begin{itemize}
    \item Leave the 'original\_study/' untouched - simply \textbf{copy over any relevant files} into the 'reproduction/' folder before running or modifying them.
    \item Use a \textbf{notebook} (.ipynb or .Rmd) when running the code (a ``literate programming approach"), as this means the code and outputs can be shared easily. The notebook can be made available to view within the Quarto site by setting it as part of the toctree in '\_quarto.yml'.
    \item Continue attempting to reproduce each item until they feel they have been \textbf{successful} (as defined in section \ref{sec:success}) - or they run out of time (from the maximum forty hours allowed).
    \item \textbf{Troubleshoot issues}, contacting the study authors if necessary (as detailed in section \ref{sec:troubleshoot})
\end{itemize}

\textbf{Remember!} Continue taking detailed notes and timings in the logbook, including each success and issue, and make note of any changes made to the model code. Whilst keeping notes in the logbook, it is recommended that \textbf{in-progress outputs are copied over} to the blog post folder (e.g. .png file for a figure), so that the researcher can easily and visually share progress within the log.

\vspace{0.5cm}
\subsubsection{Successful reproduction} \label{sec:success}

For \textbf{each item in the scope}, the researcher should decide whether it has been \textbf{successfully reproduced}. A binary decision should be made for each item (and none should be labelled as 'partial success').

This is a \textbf{subjective} decision. A successful reproduction does \textbf{not require that exactly the same results} are found. An item can be considered successfully reproduced if \textbf{minimal variation} is observed from the original results.

As an example, if it is possible to produce a table with some numbers being a match or very similar, but some numbers being substantially different, then this would be classed as having \textbf{not} been successfully reproduced. If however all aspects of the item were reproduced with reasonable similarity, this can be classed as \textbf{successful} reproduction.

Further recommendations:
\begin{itemize}
    \item \textbf{Figures} - these are compared visually. Researchers should be unconcerned by \textbf{minor differences in presentation}.
    \item \textbf{Numbers} - researchers should calculate and report the \textbf{percentage difference} in results between the original study and those obtained by the researcher. As reported by Wood et al. (2018),\autocite{wood_push_2018, wood_replication_2018} a meaningful difference in a value will vary between studies, and so it is difficult to set a single rule on what is or is not a minor difference. As such, researchers should follow a similar approach to Schwander et al. (2021),\autocite{schwander_replication_2021} considering whether the figure is reproducible at varying levels of percentage error (5\%, 10\% and 20\%). However, they \textbf{should then use their judgement to decide whether the item has been reproduced}. This is similar to one of the definitions for ``reproduction success" proposed by McManus et al. (2019)\autocite{mcmanus_can_2019} - ``\textit{Results... vary only by XX\% compared to the original, AND are consistent with the original conclusions}" - incorporating both numerical comparison and allowance for variability in whether this constitutes a meaningful difference from the original results.
\end{itemize}

Before concluding the reproduction, this decision should be run by at least one other researcher on the project, to ensure there is a \textbf{consensus decision} on whether the items were successfully reproduced. This conversation and the decision made should be included in the timing and recorded in the logbook.

In assessing reproduction success, it is important to note (as in Laurinavichyute et al. (2022)\autocite{laurinavichyute_share_2022} and Wood et al. (2018)\autocite{wood_push_2018}) that the focus is \textbf{not} on the quality or robustness of the original results, or whether the main claims of the study are consistent. Instead, the focus is on whether it was possible to reproduce the article's results \textbf{within a reasonable margin of error} (given that we do expect a little variation, since discrete-event simulations are stochastic models, and may not have been fully controlled using random seeds or with any environment differences).

It is important to \textbf{clearly timestamp} in the logbook once the decision of ``successful reproduction" has been made for each of the items in the scope.

\vspace{0.5cm}
\subsubsection{Troubleshooting} \label{sec:troubleshoot}

Researchers should \textbf{troubleshoot} any issues encountered (including \textbf{making changes to the provided code}). In allowing modification and writing of code, our intention is that researchers try \textbf{as much as possible} to attempt to reproduce from the scope. The allowance of writing new code is similar to the approach of Krafczyk et al. (2021)\autocite{krafczyk_learning_2021} and the ACRe project\autocite{berkeley_initiative_for_transparency_in_the_social_sciences_guide_2022}. Examples of changes they may need to make include:
\begin{itemize}
    \item Correcting paths to files
    \item Correcting the versions of software, or adding missing packages or libraries
    \item Fixing errors in the code
    \item Adding code to produce an item in the scope, if not otherwise provided
    \item Adding a method for \textbf{controlling randomness} in the simulation (if not otherwise set up), so that they can get consistent results with themselves between re-runs of your notebook
    \item Adding a warm-up period (if suspected but not included in the code)
\end{itemize}

Troubleshooting can include asking \textbf{advice from other members of the STARS team}. In these cases, the researcher should ensure that they include a record the time spent, everything discussed, and any recommendations made.

\vspace{0.5cm}
\subsubsection{Contacting the authors}

Despite troubleshooting, the researcher may remain unable to run the code, or have large discrepancies with the original paper. In this case, once troubleshooting is exhausted, they should \textbf{contact the original author}. This email should:
\begin{itemize}
    \item \textbf{Recap} the project (since our last email, when we informed them about the study).
    \item Link to the \textbf{Quarto site} with the documented reproduction attempt and list of issues that require resolution. Make sure the description of the problem is specific (e.g. identifying line in paper and place in code where we think something is missing, or where an issue is occurring).
    \item \textbf{Ask for suggestions} on an alternative course of action for issues, or for the complete code/data if missing.
\end{itemize}

If there is no response in two weeks, the researcher should contact them again. If there is still no response two weeks later, this can be marked as non-response. When emailing authors, it is suggested to follow the guidance on language and adapt from the \textbf{email templates} provided by ACRe in the chapter ``Guidance for Constructive Communication Between Reproducers and Original Authors" from their guide.\autocite{berkeley_initiative_for_transparency_in_the_social_sciences_guide_2022} The allowance of contacting authors is similar to the approaches of several studies,\autocite{krafczyk_learning_2021,wood_push_2018,berkeley_initiative_for_transparency_in_the_social_sciences_guide_2022,hardwicke_analytic_2021,konkol_computational_2019} with a maximum of four weeks for responses, as in Konkol et al. (2018)\autocite{konkol_computational_2019}. This approach does however differ from Laurinavichyute et al. (2022)\autocite{laurinavichyute_share_2022} who did not contact authors, since they considered reproducibility to be only about the available data and procedures and not anything shared privately.\autocite{laurinavichyute_share_2022}

\vspace{0.5cm}
\subsubsection{Running out of time}

Once forty hours has passed, the reproduction stage stops and the researcher should move onto the evaluation stage for this paper. Before moving on, they should ensure:
\begin{enumerate}
    \item If the model has no control for randomness (i.e. not implemented in original model and not yet added by researcher during troubleshooting stage) then they should select an alternative method for getting stable simulation results. This can be by doing a very large number of replications, or by assessing the required number of replications for a stable simulation.
    \item Also, they should make a final decision on the reproduction success of each item.
\end{enumerate}

\vspace{0.5cm}
\subsection{Finishing up}

\subsubsection{Tidy up notebook and create reproduction success page}

\textbf{Tidy} the reproduction notebook, so it simply produces each of the items in the scope, and clearly state how each section relates to the original article (e.g. captioning 'Reproduction attempt for Figure 2').

Using the \textbf{template page} ('evaluation/reproduction\_success.qmd'), show each item from the scope (as in the original article) alongside our best reproduction attempts (if possible under the article's license). Include the decision on the reproduction success for each item (along with any justification for this decision).
\section{Second stage: evaluation}

This section is completed \textbf{after} the attempted reproduction (so as to not interfere with timings).

\textbf{Remember!} Record progress in your \textbf{logbook} and \textbf{time spent} on each task.

\textbf{Important:} This evaluation is based on the \textbf{original} journal article or repository from the author (as in 'original\_study/'), and not on the repository that was created whilst reproducing this study ('reproduction/'). If the original study had multiple repositories to choose from (e.g. development and archived code, both prior to publication date), remember to \textbf{refer to both of them} if there are any differences between them.

\textbf{Getting a second opinion:} If the researcher is uncertain about any criteria, they should note these in the \textbf{logbook}. Any criteria that were \textbf{unmet} or \textbf{uncertain} should then be discussed with at least one other researcher on the project to get a second opinion. Record the discussion (and its timing) in the logbook, and explain and justify the choices for uncertain items.

\vspace{0.5cm}
\subsection{Badges} \label{sec:badges}

Several organisations and journals have developed badges which can be displayed alongside a research article to indicate how open and potentially reproducible it is, as detailed in Appendix \ref{appendix:badges}. These include the National Information Standards Organisation (NISO),\autocite{association_for_computing_machinery_acm_artifact_2020} the Association for Computing Machinery (ACM),\autocite{association_for_computing_machinery_acm_artifact_2020} the Institute of Electrical and Electronics Engineers (IEEE),\autocite{institute_of_electrical_and_electronics_engineers_ieee_about_nodate} the Center for Open Science (COS)\autocite{blohowiak_badges_2023}, and the journal Psychological Science.\autocite{hardwicke_transparency_2023,association_for_psychological_science_aps_psychological_2023}

We will evaluate the original study artefacts (repository) against badges that relate to code (and not those specific to data), due to the nature of DES models (where ``data" is often just parameters as part of the model script, with perhaps a few additional parameters in a separate data file within the repository). The \textbf{badges we will evaluate against} are:
\begin{itemize}
    \item ``Open objects" badges:
    \begin{itemize}
        \item NISO ``Open Research Objects (ORO)" and ``Open Research Objects - All (ORO-A)"\autocite{niso_reproducibility_badging_and_definitions_working_group_reproducibility_2021}
        \item ACM ``Artifacts Available"\autocite{association_for_computing_machinery_acm_artifact_2020}
        \item COS ``Open Code"\autocite{blohowiak_badges_2023}
        \item IEEE ``Code Available"\autocite{institute_of_electrical_and_electronics_engineers_ieee_about_nodate}
    \end{itemize}
    \item ``Object review" badges:
    \begin{itemize}
        \item ACM ``Artifacts Evaluated - Functional" and ``Artifacts Evaluated - Reusable"\autocite{association_for_computing_machinery_acm_artifact_2020}
        \item IEEE ``Code Reviewed"\autocite{institute_of_electrical_and_electronics_engineers_ieee_about_nodate}
    \end{itemize}
    \item ``Reproduced" badges:
    \begin{itemize}
        \item NISO ``Results Reproduced (ROR-R)"\autocite{niso_reproducibility_badging_and_definitions_working_group_reproducibility_2021}
        \item ACM ``Results Reproduced"\autocite{association_for_computing_machinery_acm_artifact_2020}
        \item IEEE ``Code Reproducible"\autocite{institute_of_electrical_and_electronics_engineers_ieee_about_nodate}
        \item Psychological Science ``Computational Reproducibility"\autocite{hardwicke_transparency_2023,association_for_psychological_science_aps_psychological_2023}
    \end{itemize}
\end{itemize}

The researcher should use the \textbf{provided template} (evaluation/badges.qmd)\autocite{heather_template_2024} to assess whether the artefacts from the original study meet the criteria for each of these badges. A \textbf{binary} decision is made for each criteria (as being either met or not met).

\vspace{0.5cm}
\subsection{STARS framework}

The artefacts (repository) associated with the original study will be evaluated against the \textbf{STARS framework}, which has essential and optional recommendations for sharing research artefacts from healthcare simulation studies. This framework was designed by Monks et al. (2024)\autocite{monks_towards_2024} to complement and build on general open science recommendations from the Turing Way,\autocite{the_turing_way_community_turing_2022} Taylor et al. (2017),\autocite{taylor_open_2017} and the Open Modelling Foundation (OMF) minimal and ideal reusability standards.\autocite{the_open_modeling_foundation_omf_reusability_2024}

The researcher should use the \textbf{provided template} ('evaluation/artefacts.qmd')\autocite{heather_template_2024} to assess whether the artefacts from the original study meet the recommendations from this framework. Each criteria are evaluated as being ``fully", ``partially" or ``not met".

\vspace{0.5cm}
\subsection{Reporting guidelines} \label{sec:reporting}

The \textbf{journal article} will be evaluated against two reporting guidelines for discrete-event simulation studies:
\begin{itemize}
    \item STRESS-DES: Strengthening The Reporting of Empirical Simulation Studies (Discrete-Event Simulation) - Monks et al. (2019)\autocite{monks_strengthening_2019}
    \item The generic reporting checklist for healthcare-related discrete event simulation studies derived from the International Society for Pharmacoeconomics and Outcomes Research Society for Medical Decision Making (ISPOR-SDM) Modeling Good Research Practices Task Force reports - Zhang et al. (2020)\autocite{zhang_reporting_2020}
\end{itemize}

The researcher should use the \textbf{provided template} ('evaluation/reporting.qmd')\autocite{heather_template_2024} to assess whether the criteria from these guidelines are met by the journal article (including the supplementary material, although not including the code unless the article specifically refers to it for providing particular information). Each criteria are evaluated as being ``fully", ``partially" or ``not met", with detailed evidence provided to support these claims (such as quotations from the article). If a criteria is not met by the original study, the researcher is welcome to make a suggestion in the evidence column of what they think the likely answer for that criteria might be.
\section{Summary report and research compendium}

\subsection{Computational reproducibility report} \label{sec:report}

Use the \textbf{provided template} to produce a simple summary report for the reproducibility assessment. This template should guide you to include:

\begin{itemize}
    \item Short study description
    \item Citation to original study
    \item Number and percentage of items reproduced from scope and time elapsed
    \item Required troubleshooting steps
    \item Display of reproduced items alongside original study items
    \item Percent stacked bar chart displaying the proportion of criteria met for each of the evaluations against guidelines
\end{itemize}

\vspace{1cm}
\subsection{Research compendium} \label{sec:compendium}

Once the computational reproduction has been completed, the repository will be restructured into a "\textbf{research compendium}". This a term first introduced by Gentleman and Lang 2007\autocite{gentleman_statistical_2007} which they define as "\textit{both a container for the different elements that make up the document and its computations (i.e. text, code, data,. . . ), and as a means for distributing, managing and updating the collection.}"\autocite{gentleman_statistical_2007} Marwick et al. 2018 define a research compendium as having three key components:
\begin{enumerate}
    \item Files organised according to convention
    \item Seperate data, methods and outputs
    \item Specifying the environment used for the analysis.\autocite{marwick_packaging_2018}
\end{enumerate}

A research compendium might also be referred to as a "\textbf{reproducibility file bundle}",\autocite{arguillas_10_2022} or as a "\textbf{reproduction package}".\autocite{krafczyk_learning_2021} Although not required to be structured as a package, this can be helpful in providing a structure for dependency management and file organisation, and for continuous integration of automated code testing.\autocite{marwick_packaging_2018} The repository will be modified as per the proposed structure below. Some of these modifications may already be in place from having gone through the reproduction steps above.

Our primary motivation in doing this step is to make it \textbf{easy and clear for someone to re-run} our reproduction, whilst making relatively \textbf{minimal changes} to the code itself. Hence, we are not necessarily amending to meet all of the recommendations for best practice of sharing DES models.

\vspace{0.5cm}
\subsubsection{Modify repository}

Make the following changes to the \textbf{reproduction/} folder, if not already implemented:
\begin{itemize}
    \item \textbf{Have seperate folders for data, methods and outputs} - as recommended by Marwick et al. 2018\autocite{marwick_packaging_2018}. The exception for this change is parameters coded into the scripts (since these would require a large amount of work and restructuring to seperate from the scripts, contrary to our motivation in this stage).
    \item Create \textbf{tests} which check whether a user is able to the same reproduce results as we obtained during the reproduction, based on comparison of csv files.
    \item Create \textbf{Dockerfile} and double-check it works (build image and run model notebook/s).
    \item Set-up GitHub action to publish the Docker image on the \textbf{GitHub container registry}.
    \item \textbf{Model notebook/s} should contain:
    \begin{itemize}
        \item Run time recorded within the model notebook.
        \item Clearly states which parts of the notebook produce each item from the reproduction scope.
    \end{itemize}
    \item \textbf{README} should contain:
    \begin{itemize}
        \item Original study citation.
        \item Simple summary of the model (potentially incorporating any diagrams of the model that were provided).
        \item Scope of the reproduction (including images of the figures/tables from the original study).
        \item Overview of the repository.
        \item Instructions for setting up the environment.
        \item Instructions for running the model (and reproducing items from the scope).
        \item Instructions for running the test/s.
        \item Hardware and software specs for the computer used for reproduction (machine, RAM, operating system and version).
        \item Run time for the model.
        \item Instructions for citation.
        \item Short description of license.
    \end{itemize}
    \item Ensure \textbf{Quarto site} displays the reproduction README and notebook/s.
\end{itemize}

\vspace{0.5cm}
\subsubsection{Test-run with second team member}

Once the research compendium is complete, a \textbf{second researcher} on the team should attempt to use it and confirm if they were able \textbf{reproduce} the results of the first researcher, and to check the compendium for \textbf{clarity}. This is similar to the approach of Krafczyk et al. 2021.\autocite{krafczyk_learning_2021} It should be recorded within the \textbf{logbook}.

\vspace{1cm}
\subsection{Archive on Zenodo}

Once modification of the repository is completed, a new GitHub release should be created to archive this in Zenodo (with record in CHANGELOG).

\vspace{0.5cm}
\subsection{Inform the authors}

\textbf{Email the authors} again to:
\begin{itemize}
    \item Let them know we have finished the assessment
    \item Include link to GitHub and Zenodo.
    \item Let them know how we are going to use the results from this work (i.e. lessons from reproduction, and guide framework design, and that this was not about validity of results).
\end{itemize}

% Sections after this point are Appendices
\appendix
\section{Appendix: Introductory email template} \label{appendix:email}

\textbf{Subject:} Reproducing results from [Title of paper]

Dear [Title (e.g. ``Dr")] [Last name of corresponding author],

I am contacting you about your paper titled [Title] which was published in [Journal] in [Year] (vol [Volume], no. [Number]), [Link].  

I am part of a research team working on a UKRI Medical Research Council funded project called Sharing Tools and Artefacts for Reproducible Simulations (STARS) in healthcare (\url{https://gtr.ukri.org/projects?ref=MR%2FZ503915%2F1}). We are writing to inform you that we were impressed with your paper and have selected it to be part of a STARS reproduction test.  

STARS is a methods and training study investigating open science in healthcare simulation (specifically models written in Python or R). We aim to identify best practice and learn how we can support researchers make their models and code available, robust, reusable, reproducible, and preserved long term.

As part of STARS, we identified your paper as one of the small number of simulation studies where the code and model has been made available with an open license.  The team would like to commend you for making your work open and for being a pioneer in this area.

The paper is one of several we have selected from the literature.  Our initial study in STARS will attempt to run the published code and reproduce the results on our own computers. The work will record the challenges in installing, running and reproducing the results that are reported in the selected papers.  As we are working with multiple studies our aim is to learn general lessons about sharing simulation models and move the field forward. This work is not about critiquing individual studies.  

We have a full protocol describing the work available online: [Link to protocol on Zenodo]

We may contact you, if we are unable to run your model or obtain similar results.  We appreciate that your time is valuable and we will only do so if we have exhausted all options.  It is, of course, at your discretion if you reply to our query or not.  We will also send you a copy of the results of the reproduction at the end of the study. The paper will be correctly cited and we will abide to the license conditions attached to the code.

There is nothing for you to do, and you do not need to reply to this email. However, do please feel free to contact us if you have any questions. Thank you again for making your work open so that others can learn from it.

[Reproducer full name]

[Reproducer job title], University of Exeter.

On behalf of the STARS team.
\section{Appendix: Badges} \label{appendix:badges}

A badge is an that is displayed alongside a published research article. There are currently several different badges from various organisations that can be used to indicate how open and potentially reproducible an article and its artifacts are. We have identified journal badges provided by the following organisations/journals:
\begin{itemize}
    \item National Information Standards Organisation (NISO)
    \item Association for Computing Machinery (ACM)
    \item Institute of Electrical and Electronics Engineers (IEEE)
    \item Center for Open Science (COS)
    \item Springer Nature
    \item Psychological Science
\end{itemize}

Badges related to reproducibility are typically across three categories - "open objects", "object review" and "reproduced". However, there are other badges available that are related to reproducibility. These include badges for pre-registration, like the COS "Preregistered" badge.\autocite{blohowiak_badges_2023} There are also badges for replication, which is when an independent study on the same question find consistent results (potentially with new artifacts and methods). Examples of this are the NISO "Results Replicated (RER)" badge\autocite{niso_reproducibility_badging_and_definitions_working_group_reproducibility_2021} and IEEE "Code Replicated" and "Dataset Replicated" badges.\autocite{institute_of_electrical_and_electronics_engineers_ieee_about_nodate}

In some journals, these criteria are set as requirements for publication with the journal, rather than as badges. An example of this is the Psychological Science journal, which recently transitioned from awarding COS badges to making them requirements.\autocite{hardwicke_transparency_2023}

This focuses on badges awarded by journals, but there are examples of badges that can be added to a repository if authors have followed a particular framework, either by self-allocating the badge or going through a review process. An example of this is Van Lissa et al. 2021\autocite{van_lissa_worcs_2021} who provide a package to facilitate use of a reproducible workflow in R, and suggest that a badge can be added to the README.md file of that repository if the package is used.\autocite{van_lissa_worcs_2021}

\subsection{"Open objects" badges}

"Open objects" badges relate to permanently archiving digital objects (data, code) in a public repository with a persistent identifier and open license.\autocite{niso_reproducibility_badging_and_definitions_working_group_reproducibility_2021} Examples include:
\begin{itemize}
    \item NISO "Open Research Objects (ORO)"\autocite{niso_reproducibility_badging_and_definitions_working_group_reproducibility_2021}
    \item ACM "Artifacts Available"\autocite{association_for_computing_machinery_acm_artifact_2020}
    \item COS "Open Data”, ”Open Materials" and "Open Code"\autocite{blohowiak_badges_2023}
    \item IEEE "Code Available” and ”Datasets Available"\autocite{institute_of_electrical_and_electronics_engineers_ieee_about_nodate}
    \item Springer Nature "Badge for Open Data"\autocite{springer_nature_springer_2018}
\end{itemize}

The table below compares the criteria for each of these badges (based on the sources cited above).

\begin{table}[H]
\centering
\caption{"Open objects" badge criteria}
% Space between caption and table
\vspace{0.2cm}
% Identifier for the table
\label{table:badges}
% Minimum column width
\tymin=2cm
% Add space between table rows
{\renewcommand{\arraystretch}{1.2}
    \begin{tabulary}{\linewidth}{@{}LCCCCC@{}}
          \toprule
          Criteria & NISO & ACM & COS & IEEE & Springer Nature
          \\\midrule
          Criteria apply to code & \checkmark & \checkmark & \checkmark & \checkmark & X
          \\\addlinespace
          Criteria apply to data & \checkmark & \checkmark & \checkmark & \checkmark & \checkmark
          \\\midrule
          Stored in permanent archive that is publicly and openly accessible  & \checkmark & \checkmark & \checkmark & X & \checkmark*
          \\\addlinespace
          Has a persistent identifier & \checkmark & \checkmark  & \checkmark & X & \checkmark
          \\\addlinespace
          Includes an open license describing terms of reuse & \checkmark & X & \checkmark & X & X
          \\\addlinespace
          Complete set of materials are shared (needed to fully reproduce article) & \checkmark & X & \checkmark & \checkmark & X
          \\\addlinespace
          Has metadata that sufficiently describes data/code to enable reproduction (e.g. package versions) & X & X & \checkmark & X & X
          \\\addlinespace
          Manuscript includes data availability statement & X & X & X & X & \checkmark
          \\\bottomrule
    \end{tabulary}
}
\end{table}

% Table footnotes
\vspace*{-1.5\baselineskip}
\footnotesize 
* Springer Nature Open Data badge only requires data to be in a public repository - does not require it to be permanent or openly accessible - although does require it to have a persistent identifier.

Abbreviations: ACM, Association for Computing Machinery; COS, Center for Open Science; IEEE, Institute of Electrical and Electronics Engineers; NISO, National Information Standards Organisation.
\normalsize 
\\

\subsection{"Object review" badges}

"Object review" badges relate to the digital objects (i.e. data, code) being reviewed according to the criteria of the badge issuer.\autocite{niso_reproducibility_badging_and_definitions_working_group_reproducibility_2021} Examples include:
\begin{itemize}
    \item NISO "Research Objects Reviewed (ROR)"\autocite{niso_reproducibility_badging_and_definitions_working_group_reproducibility_2021}
    \item ACM "Artifacts Evaluated - Functional" and "Artifacts Evaluated - Reusable"\autocite{association_for_computing_machinery_acm_artifact_2020}
    \item IEEE "Code Reviewed" and "Datasets Reviewed"\autocite{institute_of_electrical_and_electronics_engineers_ieee_about_nodate}
\end{itemize}

Their criteria are summarised and compared below. The NISO badge is not included as it does not have criteria, but just states that badges in this category would evaluate against a specified set of criteria.

\begin{table}[H]
\centering
\caption{"Open review" badge criteria}
% Space between caption and table
\vspace{0.2cm}
% Identifier for the table
\label{table:badges}
% Minimum column width
\tymin=2cm
% Add space between table rows
{\renewcommand{\arraystretch}{1.2}
    \begin{tabulary}{\linewidth}{@{}LCC@{}}
          \toprule
          Criteria & ACM & IEEE
          \\\midrule
          Complete code/data provided to reviewer (i.e. checking that it includes all components for producing the article outputs) & \checkmark & \checkmark
          \\\addlinespace
          Description of artefacts (minimal/sufficient, or carefully documented) & \checkmark & X
          \\\addlinespace
          Artefacts are well structured to facilitate reuse, adhering to norms and standards of research community & \checkmark & X
          \\\bottomrule
    \end{tabulary}
}
\end{table}

% Table footnotes
\vspace*{-1.5\baselineskip}
\footnotesize
Abbreviations: ACM, Association for Computing Machinery; IEEE, Institute of Electrical and Electronics Engineers;
\normalsize 
\\

\subsubsection{"Reproduced" badges}

"Reproduced" badges are awarded when an independent party regenerates the article results using author objects.\autocite{niso_reproducibility_badging_and_definitions_working_group_reproducibility_2021} Examples include:
\begin{itemize}
    \item NISO "Results Reproduced (ROR-R)"\autocite{niso_reproducibility_badging_and_definitions_working_group_reproducibility_2021}
    \item ACM "Results Reproduced"\autocite{association_for_computing_machinery_acm_artifact_2020}
    \item IEEE "Code Reproducible" and "Dataset Reproducible"\autocite{institute_of_electrical_and_electronics_engineers_ieee_about_nodate}
    \item Psychological Science "Computational Reproducibility"\autocite{hardwicke_transparency_2023,association_for_psychological_science_aps_psychological_2023}
\end{itemize}

The criteria are summarised in the table below. It should be noted that ACM specify that results are reproduced "in part" using artefacts from the author, and that exact reproduction is not required but that results should be within an acceptable range for experiments of that type.\autocite{association_for_computing_machinery_acm_artifact_2020} Whether deviation in results or any modification of the author code (such as minor troubleshooting) is permissible is not detailed within the viewed criteria for the other badges.

\begin{table}[H]
\centering
\caption{"Reproduced" badge criteria}
% Space between caption and table
\vspace{0.2cm}
% Identifier for the table
\label{table:badges}
% Minimum column width
\tymin=2cm
% Add space between table rows
{\renewcommand{\arraystretch}{1.2}
    \begin{tabulary}{\linewidth}{@{}LCCCC@{}}
          \toprule
          Criteria & NISO & ACM & IEEE & Psychological Science
          \\\midrule
          Independent party regenerated results using the authors' research artefacts & \checkmark & \checkmark & \checkmark & \checkmark
          \\\addlinespace
          Reproduced within approximately one hour (excluding compute time) & X & X & X & \checkmark
          \\\addlinespace
          Requires data and scripts to be well-organised, clearly documented and with a README file with step-by-step instructions on how to reproduce results in the manuscript & X & X & X & \checkmark
          \\\bottomrule
    \end{tabulary}
}
\end{table}

% Table footnotes
\vspace*{-1.5\baselineskip}
\footnotesize
Abbreviations: ACM, Association for Computing Machinery; IEEE, Institute of Electrical and Electronics Engineers; NISO, National Information Standards Organisation.
% \section{Appendix: STARS framework} \label{appendix:stars}

The essential and optional components of the STARS framework are presented in Tables \ref{table:stars-essential} and \ref{table:stars-optional}, as from Monks et al. (2024).\autocite{monks_towards_2024}

\hl{add citations and abbreviations}

\begin{table}[H]
\centering
\caption{Essential components of STARS framework}
% Space between caption and table
\vspace{0.2cm}
% Identifier for the table
\label{table:stars-essential}
% Minimum column width
\tymin=3cm
% Add space between table rows
{\renewcommand{\arraystretch}{1.2}
    \begin{tabulary}{\linewidth}{@{}LL@{}}
        \toprule
        Component & Description
        \\\midrule
        Open license & Free and open-source software (FOSS) license (e.g. MIT, GNU Public License (GPL))
        \\\addlinespace
        Dependency management & Specify software libraries, version numbers and sources (e.g. dependency management tools like pip virtual environment, conda environment, poetry)
        \\\addlinespace
        FOSS model & Coded in FOSS language (e.g. R, Julia, Python)
        \\\addlinespace
        Minimum documentation & Minimal instructions (e.g. in README) that overview (a) what model does, (b) how to install and run model to obtain results, and (c) how to vary parameters to run new experiments
        \\\addlinespace
        ORCID & ORCID for each study author
        \\\addlinespace
        Citation information & Instructions on how to cite the research artefact (e.g. CITATION.cff file)
        \\\addlinespace
        Remote code repository & Code available in a remote code repository (e.g. GitHub, GitLab, BitBucket)
        \\\addlinespace
        Open science archive & Code stored in an open science archive with FORCE11 compliant citation and guaranteed persistance of digital artefacts (e.g. Figshare, Zenodo, the Open Science Framework (OSF), and the Computational Modeling in the Social and Ecological Sciences Network (CoMSES Net))
        \\\bottomrule
    \end{tabulary}
}
\end{table}

\begin{table}[H]
\centering
\caption{Optional components of STARS framework}
% Space between caption and table
\vspace{0.2cm}
% Identifier for the table
\label{table:stars-optional}
% Minimum column width
\tymin=3cm
% Add space between table rows
{\renewcommand{\arraystretch}{1.2}
    \begin{tabulary}{\linewidth}{@{}LL@{}}
        \toprule
        Component & Description
        \\\midrule
        Enhanced documentation & Open and high quality documentation on how the model is implemented and works  (e.g. via notebooks and markdown files, brought together using software like Quarto and Jupyter Book). Suggested content includes:\newline• Plain english summary of project and model\newline• Clarifying license\newline• Citation instructions\newline• Contribution instructions\newline• Model installation instructions\newline• Structured code walk through of model\newline• Documentation of modelling cycle using TRACE\newline• Annotated simulation reporting guidelines\newline• Clear description of model validation including its intended purpose
        \\\addlinespace
        Documentation hosting & Host documentation (e.g. with GitHub pages, GitLab pages, BitBucket Cloud, Quarto Pub)
        \\\addlinespace
        Online coding environment & Provide an online environment where users can run and change code (e.g. BinderHub, Google Colaboratory, Deepnote)
        \\\addlinespace
        Model interface & Provide web application interface to the model so it is accessible to less technical simulation users
        \\\addlinespace
        Web app hosting & Host web app online (e.g. Streamlit Community Cloud, ShinyApps hosting)
        \\\bottomrule
    \end{tabulary}
}
\end{table}

% Print bibliography and include in table of contents
\printbibliography[heading=bibintoc]

\end{document}
